\section{Public-Code Mutation-Shape Stress Test}
\label{sec:real-code-audit}

The internally authored authority cases establish controlled Patch behavior, but they do not show how the current single-target \texttt{Change} abstraction fares against mutation patterns found in mature public software. We therefore inspect 15 commit-bound mutation sites from Excalidraw, VS Code, and JupyterLab. This is an illustrative stress test, not an integration study, repository-mining study, or estimate of ecosystem prevalence. None of these projects is translated to, ported to, or executed by Patch.

For each site we ask two separate questions: (1) can the state transition be described as one Patch-style semantic Change over one logical target using a set/add/remove interpretation, and (2) does the application itself own the commit route, or is persistence controlled by a foreign host API? A site can therefore have a simple semantic shape while still falling outside Patch governance.

\paragraph{Method.}
The corpus manifest under \path{evaluation/real-code-audit} freezes exact repository commits, source paths, and textual anchors. All three projects are from the JavaScript/TypeScript ecosystem; the intended variation is in state context, not programming-language family. The sample includes application-owned editor state, collection mutation, and host-backed persistence. Candidate files were located with state-specific searches and frozen before classification. Each site is classified along two dimensions: semantic shape and commit-route ownership. A \emph{single-target} site has a set/add/remove interpretation over one logical state target. A \emph{multi-target bundle} changes several related fields in one source update. \emph{Application-owned} state is ordinary application/editor state; \emph{host-backed} state is committed by an external service such as VS Code Memento or JupyterLab \texttt{IStateDB}.

The classifications were performed by one coder using the frozen code and the stated rubric. No inter-rater agreement statistic is reported. The corpus and rationale are published so a later replication can add independent coding. Host-backed updates are counted conservatively: a one-key Memento update may have a compact set-like transition, but Patch does not own that commit route. Multi-field updates are likewise not silently decomposed and counted as direct fits when atomic grouping may matter.

\begin{table}[t]
\centering
\caption{Mutation-shape stress test. Counts describe only the frozen 15-site corpus.}
\label{tab:real-code-audit}
\small
\begin{tabular}{@{}lrp{67mm}@{}}
\toprule
Classification & Sites & Meaning for Paper~1 \\
\midrule
Application-owned + single-target & 5 & closest direct representational fit to the current Paper~1 abstraction \\
Host-backed + single-target & 7 & simple semantic shape, but persistence crosses a foreign state API \\
Multi-target bundle & 3 & grouped or atomic multi-target semantics are outside Paper~1 \\
\midrule
All single-target shapes & 12 & set/add/remove shape irrespective of who owns the commit route \\
Application-owned state & 8 & application/editor state irrespective of mutation shape \\
\bottomrule
\end{tabular}
\end{table}

\paragraph{Observations.}
The most conservative direct-fit count is 5/15: five sites are both application-owned and single-target. Seven additional sites also have a single-target shape but commit through a foreign host-state API, so they cannot be counted as Patch-controlled mutations. The remaining three sites are multi-target bundles. The often more flattering 12/15 single-target count is therefore secondary: it describes semantic shape only, not end-to-end applicability.

The stress test exposes two concrete boundaries. First, host-backed persistence can resemble a Patch \texttt{set} while remaining outside the Patch commit route. For example, VS Code's \texttt{globalState.update(key, true)} is set-like, yet VS Code Memento performs the persistent update. JupyterLab's recents manager makes the reconstruction issue more explicit: local \texttt{splice}/\texttt{unshift} operations express remove/add semantics, while a later \texttt{IStateDB.save} persists the aggregate recents object through a foreign API. The corpus does not test whether any history or undo metadata becomes stale, so this is evidence of separated semantic and persistence steps, not of an observed consistency defect. Second, multi-field Excalidraw \texttt{setState} updates expose grouping that the current single-target model does not claim to commit atomically. These cases are useful precisely because they are not all successes.

\paragraph{Validity boundary.}
The sample is small, purposefully selected, limited to one broad language ecosystem, and classified by a single coder. It cannot support prevalence, language-family, migration, usability, or compatibility claims. The stress test does not execute third-party code through Patch, certify host adapters, measure porting effort, or extend the Lean fragment. Its role is to provide inspectable external examples that test the current abstraction and make its boundaries concrete.
